\chapter{Chapter06 — Siegenthaler Sharpness}
%%% generated by the blueprint pipeline from FABL.Chapter06.F₂Polynomials.SiegenthalerSharpness
\section{Overview}

Declarations of the Lean module \texttt{FABL.Chapter06.F₂Polynomials.SiegenthalerSharpness}.

\section{Declarations}

\begin{definition}[Decoding a binary Boolean function]
\lean{FABL.booleanFunctionF₂Decoding}
\leanok
\uses{FABL.BooleanFunction, FABL.F₂BooleanFunction, FABL.binaryCubeSignEquiv, FABL.signEncode}
\statementsource{Analysis-of-Boolean-Functions-by-Ryan-ODonnell}{
  pdf-ffbe9c7b3306-p019-b003,
  pdf-ffbe9c7b3306-p020-b001,
  pdf-ffbe9c7b3306-p148-b003,
  pdf-ffbe9c7b3306-p149-b001,
  pdf-ffbe9c7b3306-p169-b008
}
For a function $f:(\mathbb{Z}/2\mathbb{Z})^n\to\mathbb{Z}/2\mathbb{Z}$, define a
sign-valued function on sign vectors of length $n$ by first converting the sign vector
to its corresponding binary vector, applying $f$, and encoding the resulting bit as a
sign.
\end{definition}

\begin{theorem}[Encoding after decoding]
\lean{FABL.booleanFunctionF₂Encoding_booleanFunctionF₂Decoding}
\leanok
\uses{FABL.F₂BooleanFunction, FABL.binarySignEquiv, FABL.booleanFunctionF₂Decoding, FABL.booleanFunctionF₂Encoding}
\difficulty{3}
Let $n$ be a nonnegative integer and let $f:\mathbb F_2^n\to\mathbb F_2$ be a function.
Encoding the sign-valued function obtained by decoding $f$ recovers $f$ exactly.
\end{theorem}

\begin{theorem}[Real form of binary decoding]
\lean{FABL.booleanFunctionF₂Decoding_toReal}
\leanok
\uses{FABL.BooleanFunction.toReal, FABL.F₂BooleanFunction, FABL.binaryFunctionOnSignCube, FABL.booleanFunctionF₂Decoding, FABL.realSignEncodedFunction, FABL.signEncodedFunction}
\difficulty{3}
Let $n$ be a nonnegative integer and let $f:\mathbb F_2^n\to\mathbb F_2$ be a function.
Encode each bit as a sign and transport $f$ to the sign cube using the coordinatewise
bit-to-sign bijection. Then the real-valued version of the decoded sign-valued function
equals the real sign encoding of $f$ transported from the binary cube to the sign cube.
\end{theorem}

\begin{theorem}[Encoding a parity product]
\lean{FABL.booleanFunctionF₂Encoding_parityTimes}
\leanok
\uses{FABL.BooleanFunction, FABL.F₂BooleanFunction, FABL.binaryCubeSignEquiv_apply, FABL.binarySignEquiv, FABL.booleanFunctionF₂Encoding, FABL.coordinateSum, FABL.parityFunction, FABL.parityTimes, FABL.signEncode, FABL.signEncode_add, FABL.signEncode_booleanFunctionF₂Encoding, FABL.signEncode_coordinateSum}
\proofsource{Analysis-of-Boolean-Functions-by-Ryan-ODonnell}{
  pdf-ffbe9c7b3306-p023-b001,
  pdf-ffbe9c7b3306-p148-b004,
  pdf-ffbe9c7b3306-p150-b001
}
\difficulty{4}
Let $n$ be a nonnegative integer, let $S$ be a subset of its coordinates, and let $g$ be
a sign-valued Boolean function on the $n$-dimensional sign cube. The binary encoding of
the pointwise product of the parity function on $S$ with $g$ is the pointwise sum in
$\mathbb F_2$ of the coordinate sum on $S$ and the binary encoding of $g$.
\end{theorem}

\begin{theorem}[Coordinate dependence of a decoded monomial]
\lean{FABL.booleanFunctionF₂Decoding_anfMonomial_dependsOn}
\leanok
\uses{FABL.anfMonomial, FABL.binarySignEquiv, FABL.booleanFunctionF₂Decoding, FABL.signEncode}
\difficulty{3}
Let $n$ be a nonnegative integer and let $S\subseteq\{0,\ldots,n-1\}$. The sign-valued
function on $\{\pm1\}^n$ obtained by decoding the $\mathbb{F}_2$ monomial $\prod_{i\in
S}z_i$ depends only on the coordinates in $S$: whenever two sign vectors agree on every
coordinate in $S$, the function has the same value on both.
\end{theorem}

\begin{theorem}[Card le function algebraic degree of anf coeff ne zero]
\lean{FABL.card_le_functionAlgebraicDegree_of_anfCoeff_ne_zero}
\statementsource{Analysis-of-Boolean-Functions-by-Ryan-ODonnell}{
  pdf-ffbe9c7b3306-p149-b001,
  pdf-ffbe9c7b3306-p149-b002,
  pdf-ffbe9c7b3306-p149-b004
}
\label{FABL.card_le_functionAlgebraicDegree_of_anfCoeff_ne_zero}
\leanok
\uses{FABL.F₂BooleanFunction, FABL.algebraicDegree, FABL.anfCoeff, FABL.functionAlgebraicDegree, FABL.mem_anfSupport}
\difficulty{3}
% dataset-pipeline: stub -- prose pending build_informal_statements.py
Statement pending regeneration.
\end{theorem}

\begin{theorem}[Coefficient of a sum and a monomial]
\lean{FABL.anfCoeff_coordinateSum_add_anfMonomial_of_disjoint}
\label{FABL.anfCoeff_coordinateSum_add_anfMonomial_of_disjoint}
\leanok
\uses{FABL.F₂BooleanFunction, FABL.affineCoefficients, FABL.affineFunction, FABL.anfCoeff, FABL.anfCoeff_add, FABL.anfCoeff_affineFunction, FABL.anfCoeff_anfMonomial, FABL.anfMonomial, FABL.coordinateSum, FABL.f₂CubeOfFinset, FABL.f₂CubeOfFinset_apply, FABL.f₂DotProduct_f₂CubeOfFinset}
\difficulty{5}
Let $n$ be a natural number, and let $S,T\subseteq\{0,\ldots,n-1\}$ be finite sets with
$T$ nonempty and $S\cap T=\varnothing$. The algebraic normal form coefficient indexed by
$T$ of the function that adds the sum of coordinates in $S$ to the monomial given by the
product of coordinates in $T$, over the field with two elements, is $1$.
\end{theorem}

\begin{definition}[Resilient sharpness prefix]
\lean{FABL.resilientSharpnessPrefix}
\statementsource{Analysis-of-Boolean-Functions-by-Ryan-ODonnell}{
  pdf-ffbe9c7b3306-p165-b005,
  pdf-ffbe9c7b3306-p170-b002,
  pdf-ffbe9c7b3306-p170-b004
}
\label{FABL.resilientSharpnessPrefix}
\leanok
Let $n,k$ be natural numbers with $k<n-1$. Define the prefix to consist of the first
$k+1$ coordinates, namely the coordinates indexed by $0,1,\ldots,k$.
\end{definition}

\begin{definition}[Resilient sharpness tail]
\lean{FABL.resilientSharpnessTail}
\statementsource{Analysis-of-Boolean-Functions-by-Ryan-ODonnell}{
  pdf-ffbe9c7b3306-p165-b005,
  pdf-ffbe9c7b3306-p170-b002,
  pdf-ffbe9c7b3306-p170-b004
}
\label{FABL.resilientSharpnessTail}
\leanok
\uses{FABL.resilientSharpnessPrefix}
Let $n,k$ be natural numbers with $k<n-1$. Define the tail to be the complement, among
the $n$ coordinates, of the prefix consisting of indices $0,1,\ldots,k$.
\end{definition}

\begin{theorem}[Size of the resilient sharpness prefix]
\lean{FABL.card_resilientSharpnessPrefix}
\label{FABL.card_resilientSharpnessPrefix}
\leanok
\uses{FABL.resilientSharpnessPrefix}
\difficulty{1}
Let $n$ and $k$ be natural numbers with $k<n-1$. The prefix consisting of the first
$k+1$ indices among the $n$ coordinates has cardinality $k+1$.
\end{theorem}

\begin{theorem}[Size of the resilient sharpness tail]
\lean{FABL.card_resilientSharpnessTail}
\label{FABL.card_resilientSharpnessTail}
\leanok
\uses{FABL.card_resilientSharpnessPrefix, FABL.resilientSharpnessTail}
\difficulty{3}
Let $n$ and $k$ be natural numbers with $k<n-1$. The complement, among the $n$
coordinate indices, of the prefix consisting of the first $k+1$ indices has cardinality
$n-k-1$.
\end{theorem}

\begin{definition}[Resilient degree sharpness function]
\lean{FABL.resilientDegreeSharpnessFunction}
\statementsource{Analysis-of-Boolean-Functions-by-Ryan-ODonnell}{
  pdf-ffbe9c7b3306-p165-b005,
  pdf-ffbe9c7b3306-p170-b002,
  pdf-ffbe9c7b3306-p170-b004
}
\label{FABL.resilientDegreeSharpnessFunction}
\leanok
\uses{FABL.BooleanFunction, FABL.anfMonomial, FABL.booleanFunctionF₂Decoding, FABL.parityTimes, FABL.resilientSharpnessPrefix, FABL.resilientSharpnessTail}
Let $n,k$ be natural numbers with $k<n-1$. Define a sign-valued function on sign vectors
of length $n$ by taking the parity on the prefix of coordinates $0,1,\ldots,k$ and
multiplying it by the sign encoding of the binary monomial supported on the
complementary tail coordinates.
\end{definition}

\begin{theorem}[Resilient sharpness function]
\lean{FABL.resilientDegreeSharpnessFunction_isResilient}
\label{FABL.resilientDegreeSharpnessFunction_isResilient}
\leanok
\uses{FABL.IsResilient, FABL.booleanFunctionF₂Decoding_anfMonomial_dependsOn, FABL.card_resilientSharpnessPrefix, FABL.parityTimes_isResilient_of_dependsOn_compl, FABL.resilientDegreeSharpnessFunction, FABL.resilientSharpnessTail}
\difficulty{3}
For natural numbers $n$ and $k$ with $k<n-1$, the function constructed from the first
$k+1$ coordinates and the remaining coordinates as in the resilient sharpness
construction is $k$-resilient: its nonempty Fourier coefficients on sets of at most $k$
coordinates vanish, and its mean is zero.
\end{theorem}

\begin{theorem}[Encoding the sharpness function]
\lean{FABL.booleanFunctionF₂Encoding_resilientDegreeSharpnessFunction}
\leanok
\uses{FABL.F₂BooleanFunction, FABL.anfMonomial, FABL.booleanFunctionF₂Encoding, FABL.booleanFunctionF₂Encoding_booleanFunctionF₂Decoding, FABL.booleanFunctionF₂Encoding_parityTimes, FABL.coordinateSum, FABL.resilientDegreeSharpnessFunction, FABL.resilientSharpnessPrefix, FABL.resilientSharpnessTail}
\difficulty{1}
Let $n$ and $k$ be nonnegative integers with $k<n-1$. Let
$P=\{i\in\{0,\ldots,n-1\}:i<k+1\}$ and let $T$ be its complement among the $n$
coordinates. The binary encoding of the sign function formed by multiplying the parity
on $P$ by the decoded monomial on $T$ equals the sum in $\mathbb F_2$ of the coordinate
sum on $P$ and the monomial on $T$.
\end{theorem}

\begin{theorem}[Function algebraic degree resilient degree sharpness function]
\lean{FABL.functionAlgebraicDegree_resilientDegreeSharpnessFunction}
\statementsource{Analysis-of-Boolean-Functions-by-Ryan-ODonnell}{
  pdf-ffbe9c7b3306-p149-b004,
  pdf-ffbe9c7b3306-p150-b004,
  pdf-ffbe9c7b3306-p165-b005,
  pdf-ffbe9c7b3306-p170-b004
}
\label{FABL.functionAlgebraicDegree_resilientDegreeSharpnessFunction}
\leanok
\uses{FABL.F₂BooleanFunction, FABL.anfCoeff, FABL.anfCoeff_coordinateSum_add_anfMonomial_of_disjoint, FABL.anfMonomial, FABL.booleanFunctionF₂Encoding, FABL.booleanFunctionF₂Encoding_resilientDegreeSharpnessFunction, FABL.card_le_functionAlgebraicDegree_of_anfCoeff_ne_zero, FABL.card_resilientSharpnessTail, FABL.coordinateSum, FABL.functionAlgebraicDegree, FABL.functionAlgebraicDegree_booleanFunctionF₂Encoding_le_of_isResilient, FABL.resilientDegreeSharpnessFunction, FABL.resilientDegreeSharpnessFunction_isResilient, FABL.resilientSharpnessPrefix, FABL.resilientSharpnessTail}
\difficulty{5}
% dataset-pipeline: stub -- prose pending build_informal_statements.py
Statement pending regeneration.
\end{theorem}

\begin{theorem}[Exists is resilient function algebraic degree eq sub sub one]
\lean{FABL.exists_isResilient_functionAlgebraicDegree_eq_sub_sub_one}
\statementsource{Analysis-of-Boolean-Functions-by-Ryan-ODonnell}{pdf-ffbe9c7b3306-p150-b004}
\label{FABL.exists_isResilient_functionAlgebraicDegree_eq_sub_sub_one}
\leanok
\uses{FABL.BooleanFunction, FABL.IsResilient, FABL.booleanFunctionF₂Encoding, FABL.functionAlgebraicDegree, FABL.functionAlgebraicDegree_resilientDegreeSharpnessFunction, FABL.resilientDegreeSharpnessFunction, FABL.resilientDegreeSharpnessFunction_isResilient}
\difficulty{1}
% dataset-pipeline: stub -- prose pending build_informal_statements.py
Statement pending regeneration.
\end{theorem}

\begin{definition}[Indicator of the two constant vectors]
\lean{FABL.constantPairF₂Indicator}
\leanok
\uses{FABL.F₂BooleanFunction}
For each natural number $n$, define a function from $(\mathbb{F}_2)^n$ to $\mathbb{F}_2$
that takes the value $1$ at the zero vector and at the all-ones vector, and takes the
value $0$ elsewhere.
\end{definition}

\begin{definition}[Decoded constant-pair Boolean function]
\lean{FABL.constantPairBooleanFunction}
\label{FABL.constantPairBooleanFunction}
\leanok
\uses{FABL.BooleanFunction, FABL.booleanFunctionF₂Decoding, FABL.constantPairF₂Indicator}
For each natural number $n$, the Boolean function is obtained by decoding the indicator
of the pair consisting of the zero vector and the all-ones vector in $(\mathbb{F}_2)^n$.
\end{definition}

\begin{theorem}[Characterization of the constant pair subspace]
\lean{FABL.mem_constantPairSubspace_iff}
\label{FABL.mem_constantPairSubspace_iff}
\leanok
\uses{FABL.F₂Cube, FABL.ProbabilityDensity.constantPairSubspace}
\difficulty{1}
Let $x$ be a vector in $\mathbb{F}_2^n$. It belongs to the constant pair subspace
exactly when it is the zero vector or the all-ones vector.
\end{theorem}

\begin{theorem}[Sign encoding of the constant-pair indicator]
\lean{FABL.realSignEncodedFunction_constantPairF₂Indicator}
\leanok
\uses{FABL.F₂Cube, FABL.ProbabilityDensity.constantPairSubspace, FABL.constantPairF₂Indicator, FABL.mem_constantPairSubspace_iff, FABL.realSignEncodedFunction, FABL.setIndicator, FABL.signEncodedFunction}
\difficulty{4}
For every nonnegative integer $n$, the real sign encoding of the function on $\mathbb
F_2^n$ that is $1$ at $0$ and at the all-ones vector, and $0$ elsewhere, is the function
$x\mapsto 1-2\mathbf{1}_{\{0,\mathbf{1}\}}(x)$, where $\mathbf{1}_{\{0,\mathbf{1}\}}$ is
the real indicator of the subspace consisting of those two vectors.
\end{theorem}

\begin{theorem}[Real Form of the Constant-Pair Function]
\lean{FABL.constantPairBooleanFunction_toReal}
\label{FABL.constantPairBooleanFunction_toReal}
\leanok
\uses{FABL.BooleanFunction.toReal, FABL.binaryFunctionOnSignCube, FABL.booleanFunctionF₂Decoding_toReal, FABL.constantPairBooleanFunction, FABL.constantPairF₂Indicator, FABL.realSignEncodedFunction}
\difficulty{1}
For every natural number $n$, the real-valued version of the sign-valued function
obtained from the indicator of the zero and all-ones vectors in the binary cube equals,
at each sign vector, the real sign encoding of that indicator evaluated at the
corresponding binary vector under the coordinatewise binary-to-sign identification.
\end{theorem}

\begin{theorem}[Singleton indicator is not orthogonal to constants]
\lean{FABL.f₂CubeOfFinset_not_mem_perpendicular_constantPairSubspace_of_card_one}
\leanok
\uses{FABL.F₂Cube, FABL.ProbabilityDensity.constantPairSubspace, FABL.coordinateSum, FABL.f₂CubeEquivFinset, FABL.f₂CubeOfFinset, FABL.f₂DotProduct_eq_coordinateSum_f₂Support, FABL.f₂Support, FABL.mem_constantPairSubspace_iff, FABL.mem_perpendicularSubspace_iff, FABL.perpendicularSubspace}
\difficulty{4}
For every natural number $n$ and every finite set $T\subseteq\mathrm{Fin}(n)$ with
$|T|=1$, its binary indicator vector is not orthogonal, under the coordinate dot product
over $\mathbb{F}_2$, to the subspace consisting of the zero vector and the all-ones
vector.
\end{theorem}

\begin{theorem}[Vanishing coefficient for a constant-pair indicator]
\lean{FABL.vectorFourierCoeff_realSignEncodedFunction_constantPairF₂Indicator_eq_zero}
\leanok
\uses{FABL.F₂Cube, FABL.ProbabilityDensity.constantPairSubspace, FABL.constantPairF₂Indicator, FABL.expect_vectorWalshCharacter, FABL.f₂CubeOfFinset, FABL.f₂CubeOfFinset_apply, FABL.f₂CubeOfFinset_not_mem_perpendicular_constantPairSubspace_of_card_one, FABL.perpendicularSubspace, FABL.realSignEncodedFunction, FABL.realSignEncodedFunction_constantPairF₂Indicator, FABL.setIndicator, FABL.vectorFourierCoeff, FABL.vectorFourierCoeff_eq_expect, FABL.vectorFourierCoeff_setIndicator_submodule_of_not_mem, FABL.vectorWalshCharacter}
\difficulty{5}
Let $T$ be a subset of the $n$ coordinate indices with exactly one element. The vector
Fourier coefficient at the binary vector supported on $T$ of the real sign encoding of
the function that equals $1$ at the zero vector and the all-ones vector and equals $0$
elsewhere is zero.
\end{theorem}

\begin{theorem}[Correlation Immunity of a Constant-Pair Function]
\lean{FABL.constantPairBooleanFunction_isCorrelationImmune}
\label{FABL.constantPairBooleanFunction_isCorrelationImmune}
\leanok
\uses{FABL.IsCorrelationImmune, FABL.constantPairBooleanFunction, FABL.constantPairBooleanFunction_toReal, FABL.f₂CubeEquivFinset, FABL.f₂CubeOfFinset, FABL.f₂Support, FABL.vectorFourierCoeff_eq_fourierCoeff_binaryFunctionOnSignCube, FABL.vectorFourierCoeff_realSignEncodedFunction_constantPairF₂Indicator_eq_zero}
\difficulty{4}
For every natural number $n$, define a function on the binary cube to equal $1$ at the
zero vector and the all-ones vector, and $0$ elsewhere; convert its outputs to signs
using $b\mapsto(-1)^b$ and identify the binary cube with the sign cube coordinatewise.
The resulting sign-valued function is correlation immune of order $1$: every nonempty
Fourier coefficient indexed by a set of at most one coordinate vanishes.
\end{theorem}

\begin{theorem}[Coefficient of a punctured full set]
\lean{FABL.anfCoeff_constantPairF₂Indicator_erase}
\leanok
\uses{FABL.F₂Cube, FABL.anfCoeff, FABL.constantPairF₂Indicator, FABL.f₂CubeOfFinset, FABL.f₂CubeOfFinset_apply}
\difficulty{5}
Let $n$ be a natural number and let $i$ be a coordinate in $br^n$; over the field with
two elements, consider the Boolean function that equals $1$ exactly at the zero vector
and the all-ones vector. Its algebraic normal form coefficient indexed by all
coordinates except $i$ is $1$.
\end{theorem}

\begin{theorem}[Degree of a two-point indicator]
\lean{FABL.functionAlgebraicDegree_constantPairBooleanFunction}
\label{FABL.functionAlgebraicDegree_constantPairBooleanFunction}
\leanok
\uses{FABL.anfCoeff, FABL.anfCoeff_constantPairF₂Indicator_erase, FABL.booleanFunctionF₂Encoding, FABL.booleanFunctionF₂Encoding_booleanFunctionF₂Decoding, FABL.card_le_functionAlgebraicDegree_of_anfCoeff_ne_zero, FABL.constantPairBooleanFunction, FABL.constantPairBooleanFunction_isCorrelationImmune, FABL.constantPairF₂Indicator, FABL.functionAlgebraicDegree, FABL.functionAlgebraicDegree_booleanFunctionF₂Encoding_le_of_isCorrelationImmune}
\difficulty{4}
Let $n$ be a natural number with $3\le n$. Consider the function on the Boolean cube
over the two-element field that is $1$ at the zero vector and the all-ones vector, and
$0$ elsewhere. Its algebraic degree is $n-1$.
\end{theorem}

\begin{theorem}[Correlation immunity and algebraic degree]
\lean{FABL.exists_firstOrderCorrelationImmune_functionAlgebraicDegree_eq_sub_one}
\label{FABL.exists_firstOrderCorrelationImmune_functionAlgebraicDegree_eq_sub_one}
\leanok
\uses{FABL.BooleanFunction, FABL.IsCorrelationImmune, FABL.booleanFunctionF₂Encoding, FABL.constantPairBooleanFunction, FABL.constantPairBooleanFunction_isCorrelationImmune, FABL.functionAlgebraicDegree, FABL.functionAlgebraicDegree_constantPairBooleanFunction}
\difficulty{1}
For every integer $n\ge 3$, there is a Boolean function on the $n$-dimensional sign cube
whose Fourier coefficients vanish on every nonempty set of at most one coordinate, and
whose encoding as a function on the binary cube has algebraic degree $n-1$.
\end{theorem}

\begin{theorem}[Unbalanced correlation-immune Boolean function]
\lean{FABL.exists_not_isBalanced_isCorrelationImmune_two_mul_div_three_sub_one}
\proofsource{Analysis-of-Boolean-Functions-by-Ryan-ODonnell}{
  pdf-ffbe9c7b3306-p165-b005,
  pdf-ffbe9c7b3306-p147-b006,
  pdf-ffbe9c7b3306-p151-b002
}
\label{FABL.exists_not_isBalanced_isCorrelationImmune_two_mul_div_three_sub_one}
\leanok
\uses{FABL.BooleanFunction, FABL.BooleanFunction.toReal, FABL.IsBalanced, FABL.IsCorrelationImmune, FABL.Sign, FABL.correlationImmuneAndExample, FABL.correlationImmuneAndExample_isCorrelationImmune, FABL.mean, FABL.mean_correlationImmuneAndExample}
\difficulty{4}
For every natural number $n$ divisible by $3$, there is a function from the sign cube
$\\{−1,1\\}^n$ to $\\{−1,1\\}$ whose uniform mean is nonzero and whose nonempty Fourier
coefficients of degrees at most $2n/3-1$ all vanish.
\end{theorem}
